\documentclass[12pt]{article}
\usepackage[margin=1in]{geometry}
\usepackage{amsmath,booktabs,graphicx,makecell,threeparttable}

\begin{document}

\providecommand{\StructuralRobustnessPath}{appendix/structural-robustness}

\subsection{Population Weighting and Resource-Cost Break-Even Values}
\label{app:policy-population-cost}

The benchmark allocation gives each community equal weight. We repeat the
exercise using the census adult population of each of the 144 communities.
For every posterior draw, the planner must preserve the predicted number of
adults treated under the experimental allocation and Control observability.
Thus a large and a small community no longer make equal contributions to the
coverage constraint. Both regimes use the same draw-specific target and the
same 3.5 km candidate-site set.

Table~\ref{tab:policy-population-allocations} reports the paired results. Under
equal village weights, Bracelet observability saves a median of 27 sites, with
a 95 percent credible interval of $[17,35]$. Population weighting changes the
levels only modestly: Control requires 101 sites and Bracelet 77 at the
posterior median, for a saving of 25 sites and a 95 percent interval of
$[16,31]$. Bracelet saves sites in every retained posterior draw. Mean
community--site distance rises because the allocation uses the Bracelet's
more persistent observability to consolidate distribution.
Under the preferred cluster weighted-likelihood bootstrap, the corresponding
population-weighted saving is 24 sites, with a 95 percent percentile interval
of $[7,31]$; every weighted refit saves at least one site.

\begin{table}[htbp]
  \centering
  \scriptsize
  \caption{Population weighting in the policy allocation}
  \label{tab:policy-population-allocations}
  \resizebox{\textwidth}{!}{%
    \begin{tabular}{llrrrrrr}
\toprule
Inference & Weighting & Control sites & Bracelet sites & Sites saved & $\Pr(\text{saved}>0)$ & Control km & Bracelet km \\
\midrule
Baseline posterior & Equal village weights, 3.5 km & \makecell[c]{106\\{}[105, 107]} & \makecell[c]{79\\{}[71, 90]} & \makecell[c]{27\\{}[17, 35]} & 1.00 & \makecell[c]{1.25\\{}[1.23, 1.28]} & \makecell[c]{1.93\\{}[1.64, 2.46]} \\
 & Population weighted, 3.5 km & \makecell[c]{101\\{}[101, 102]} & \makecell[c]{77\\{}[71, 86]} & \makecell[c]{25\\{}[16, 31]} & 1.00 & \makecell[c]{1.25\\{}[1.22, 1.26]} & \makecell[c]{1.92\\{}[1.63, 2.45]} \\
\addlinespace
Cluster weighted-likelihood bootstrap & Equal village weights, 3.5 km & \makecell[c]{106\\{}[104, 108]} & \makecell[c]{80\\{}[71, 100]} & \makecell[c]{26\\{}[7, 36]} & 1.00 & \makecell[c]{1.25\\{}[1.20, 1.31]} & \makecell[c]{1.90\\{}[1.40, 2.57]} \\
 & Population weighted, 3.5 km & \makecell[c]{101\\{}[99, 103]} & \makecell[c]{78\\{}[71, 96]} & \makecell[c]{24\\{}[7, 31]} & 1.00 & \makecell[c]{1.25\\{}[1.20, 1.30]} & \makecell[c]{1.91\\{}[1.39, 2.64]} \\
\bottomrule
\end{tabular}
%
  }
  \begin{minipage}{0.96\linewidth}
    \vspace{0.35em}\footnotesize
    \textit{Notes:} Baseline cells report posterior medians and 95 percent
    credible intervals. Cluster weighted-likelihood-bootstrap rows, when
    shown, report medians and percentile intervals across exponential-weighted
    modes. The coverage target is recomputed within each draw or mode under
    the weighting scheme shown. ``Control km'' and
    ``Bracelet km'' are adult-population-weighted mean one-way distances.
    The candidate set contains all 1,451 documented schools and excludes
    village--site links longer than 3.5 km.
  \end{minipage}
\end{table}

The benchmark counterfactual holds each treatment's estimated observability
schedule fixed when communities are consolidated. To assess that assumption,
we identify each community's experimental PoT with its unique nearest
candidate school. For assignments to any other candidate, we move Bracelet
observability either 50 or 100 percent of the way toward the Control schedule,
while leaving the community's own experimental-PoT observability unchanged.
This is a deliberately pessimistic sensitivity calculation rather than an
endogenous model of traffic or information flows.

Table~\ref{tab:policy-pooling-observability} shows that 50 percent attenuation
reduces the median site saving from 24 to 10, with a 95 percent interval of
$[3,20]$ and a positive saving in 99.7 percent of weighted refits. If all
away-from-experimental-PoT assignments instead receive Control observability,
the median saving falls to 3 sites, the interval is $[0,9]$, and the saving is
positive in 89.6 percent of refits. Thus moderate information dilution still
supports consolidation, whereas the full-Control bound leaves only modest and
imprecise savings.

\begin{table}[htbp]
  \centering
  \scriptsize
  \caption{Policy sensitivity to observability under site consolidation}
  \label{tab:policy-pooling-observability}
  \resizebox{\textwidth}{!}{%
    \begin{tabular}{llrrrr}
\toprule
Inference & Consolidation observability & Control sites & Bracelet sites & Sites saved & $\Pr(\text{saved}>0)$ \\
\midrule
Cluster weighted-likelihood bootstrap & No attenuation & \makecell[c]{101\\{}[99, 103]} & \makecell[c]{78\\{}[71, 96]} & \makecell[c]{24\\{}[7, 31]} & 1.00 \\
 & 50\% attenuation & \makecell[c]{101\\{}[99, 103]} & \makecell[c]{91\\{}[81, 99]} & \makecell[c]{10\\{}[3, 20]} & 1.00 \\
 & Control observability & \makecell[c]{101\\{}[99, 103]} & \makecell[c]{99\\{}[93, 103]} & \makecell[c]{3\\{}[0, 9]} & 0.90 \\
\bottomrule
\end{tabular}
%
  }
  \begin{minipage}{0.96\linewidth}
    \vspace{0.35em}\footnotesize
    \textit{Notes:} Entries are medians and 95 percent percentile intervals
    across 999 exponential cluster-weighted modes. ``Attenuation'' applies
    only when a community is assigned away from its experimental PoT and
    linearly moves Bracelet's social-image return toward Control's return.
    Targets and objective weights use census adult population.
  \end{minipage}
\end{table}

Site consolidation is not the same as resource-cost savings. We therefore
apply a transparent accounting exercise to each draw's population-weighted,
coverage-preserving allocation. Let $S_z$ be the number of funded sites,
$Q_z$ expected participants, and $K_z$ expected round-trip
participant-kilometres in observability regime $z$. Total resource cost is
\begin{equation}
  C_z=C_F S_z + \mathbf{1}\{z=B\}(0.20)Q_z + C_TK_z,
  \label{eq:policy-resource-cost}
\end{equation}
where $C_F$ is the fixed cost of operating a PoT and $C_T$ is the household
travel value per round-trip participant-kilometre. The \$0.20 Bracelet cost is
documented; $C_F$ and $C_T$ are scenario values rather than estimated welfare
parameters. Conditional on Bracelet saving sites, the break-even fixed cost is
\begin{equation}
  C_F^*=\frac{0.20Q_B+C_T(K_B-K_C)}{S_C-S_B}.
  \label{eq:policy-break-even}
\end{equation}

Table~\ref{tab:policy-break-even} and
Figure~\ref{fig:policy-break-even} report this threshold. With household
travel unpriced, the median fixed-site break-even value is \$261, with a 95
percent credible interval of $[\$205,\$414]$. At \$0.10 per round-trip
participant-kilometre, the median threshold rises to \$483
($[\$438,\$611]$). At that travel value, a \$500 fixed site cost makes the
Bracelet allocation resource-saving in 66 percent of posterior draws, while a
\$1,000 fixed cost does so in essentially every draw. The exercise therefore
supports a conditional claim: public observability reduces the number of
sites, but whether it reduces total resource costs depends on site operating
costs and the value placed on additional household travel.
The cluster weighted-likelihood results yield the same qualitative
conclusion with wider tails. At zero travel value, the median break-even site
cost is \$270, with a percentile interval of $[\$195,\$940]$. At \$0.10 per
round-trip participant-kilometre it is \$519
($[\$420,\$1{,}126]$); a \$500 site cost is resource-saving in 41 percent of
weighted refits and a \$1,000 site cost in 96 percent.

\begin{table}[htbp]
  \centering
  \scriptsize
  \caption{Break-even resource costs for the Bracelet allocation}
  \label{tab:policy-break-even}
  \resizebox{\textwidth}{!}{%
    \begin{tabular}{lrrrrrrr}
\toprule
Inference & Travel value & Break-even site cost & $\Pr(S_B<S_C)$ & $\Pr(\Delta C>0)$ at \$100 & at \$250 & at \$500 & at \$1,000 \\
\midrule
Baseline posterior & \$0.00 & \makecell[c]{\$261\\{}[\$205, \$414]} & 1.00 & 0.00 & 0.41 & 1.00 & 1.00 \\
 & \$0.05 & \makecell[c]{\$370\\{}[\$325, \$514]} & 1.00 & 0.00 & 0.00 & 0.97 & 1.00 \\
 & \$0.10 & \makecell[c]{\$483\\{}[\$438, \$611]} & 1.00 & 0.00 & 0.00 & 0.66 & 1.00 \\
 & \$0.15 & \makecell[c]{\$593\\{}[\$544, \$736]} & 1.00 & 0.00 & 0.00 & 0.00 & 1.00 \\
 & \$0.20 & \makecell[c]{\$704\\{}[\$646, \$876]} & 1.00 & 0.00 & 0.00 & 0.00 & 1.00 \\
 & \$0.25 & \makecell[c]{\$813\\{}[\$746, \$1,028]} & 1.00 & 0.00 & 0.00 & 0.00 & 0.97 \\
\addlinespace
Cluster weighted-likelihood bootstrap & \$0.00 & \makecell[c]{\$270\\{}[\$195, \$940]} & 1.00 & 0.00 & 0.42 & 0.88 & 0.97 \\
 & \$0.05 & \makecell[c]{\$392\\{}[\$314, \$1,022]} & 1.00 & 0.00 & 0.00 & 0.79 & 0.97 \\
 & \$0.10 & \makecell[c]{\$519\\{}[\$420, \$1,126]} & 1.00 & 0.00 & 0.00 & 0.41 & 0.96 \\
 & \$0.15 & \makecell[c]{\$635\\{}[\$519, \$1,243]} & 1.00 & 0.00 & 0.00 & 0.01 & 0.95 \\
 & \$0.20 & \makecell[c]{\$749\\{}[\$613, \$1,344]} & 1.00 & 0.00 & 0.00 & 0.00 & 0.90 \\
 & \$0.25 & \makecell[c]{\$862\\{}[\$706, \$1,455]} & 1.00 & 0.00 & 0.00 & 0.00 & 0.78 \\
\bottomrule
\end{tabular}
%
  }
  \begin{minipage}{0.96\linewidth}
    \vspace{0.35em}\footnotesize
    \textit{Notes:} Baseline break-even cells report the posterior median and
    95 percent credible interval for $C_F^*$ in equation
    \eqref{eq:policy-break-even}; cluster weighted-likelihood-bootstrap rows,
    when shown, use medians and percentile intervals across weighted modes.
    Travel values are dollars per round-trip
    participant-kilometre. $\Pr(\Delta C>0)$ is the posterior probability that
    Control costs exceed Bracelet costs at the stated fixed cost. Cost
    accounting is applied after minimizing sites; it is not a complete Ramsey
    welfare exercise.
  \end{minipage}
\end{table}

\begin{figure}[htbp]
  \centering
  \includegraphics[width=0.78\textwidth]{%
    \StructuralRobustnessPath/figures/policy-break-even-site-cost.pdf}
  \caption{Fixed-site cost required for Bracelet to break even}
  \label{fig:policy-break-even}
  \begin{minipage}{0.88\linewidth}
    \vspace{0.35em}\footnotesize
    \textit{Notes:} Lines are median break-even fixed costs and shaded regions
    are 95 percent posterior credible or weighted-mode percentile intervals,
    as indicated by the legend. Each point uses the population-weighted,
    coverage-preserving allocations under Control and Bracelet observability.
  \end{minipage}
\end{figure}

\end{document}
